% SGA 3, Expose III, Loop 1 English reconstruction.
% Scope: Remark 4.8.0 through the 4.10 suppression notice, stopping
% immediately before 4.11.
% Authority: Polo--Gille Exp3-14oct24.pdf, component-local pp. 51--55
% (combined-reader pp. 161--165; original side-pages 133--136), with only
% the 4.8.0--4.10 portions of the first and last pages admitted. Force-OCR
% is a locator only. English wording was compared with Jacob C. Reinhold's
% CC BY 4.0 Markdown lineage at commit
% e7a259f3f8608ad3edf9bf6eead3fd504dd2d23e.
% The Polo--Gille PDF controls prose, mathematics, numbering, notes, and the
% four source-derived Loop-1 diagram images included below.

\setcounter{footnote}{88}

% Source page: Exp3 local p. 51; combined p. 161; original side-page 133.
\medskip
\noindent\SGAInlineRefTarget{sga3:iii:remark:4-8-0}\textbf{注 4.8.0.}\footnote{编者注: 此注用于下面的命题, 现插入此处; 它原见于
原版\SGARefLink{sga3:iii:paragraph:4-10}{4.10}.}
\textit{设 \(S\) 为概形, \(S'\) 为 \(S\) 的闭子概形,
\(X\) 为 \(S\)-概形, \(Y\) 为 \(X\) 的 \(S\)-子概形,
并令 \(X'=X\times_S S'\) 及 \(Y'=Y\times_S S'\).
则存在 \(\mathscr O_{Y'}\)-模的满态射}
\[
 \mathscr N_{Y/X}\otimes_{\mathscr O_Y}\mathscr O_{Y'}
 \twoheadrightarrow \mathscr N_{Y'/X'}.
\]

事实上, 将 \(X\) 换为适当的开子概形后,
可假设 \(Y\) 是闭的, 由 \(\mathscr O_X\) 的理想
\(\mathscr I_Y\) 定义. 于是 \(\mathscr I_Y\) 在
\(\mathscr O_{X'}\) 中的像是定义 \(Y'\) 的理想 \(\mathscr I_{Y'}\),
且存在 \(\mathscr O_{Y'}\)-模的满态射
\[
 \pi:
 (\mathscr I_Y/\mathscr I_Y^2)
 \otimes_{\mathscr O_Y}\mathscr O_{Y'}
 \twoheadrightarrow
 \mathscr I_{Y'}/\mathscr I_{Y'}^2.
\]
再假设 \(\mathscr O_Y=\mathscr O_X/\mathscr I_Y\)
在 \(\mathscr O_S\) 上平坦. 则自然态射
\[
 \mathscr I_Y\otimes_{\mathscr O_X}\mathscr O_{X'}
 \longrightarrow \mathscr I_{Y'}
\]
为双射 (参见 EGA IV\textsubscript{2}, 2.1.8).

% Source page: Exp3 local p. 52; combined p. 162; original side-page 133.
因此有如下各行正合的交换图:

% Raster stand-in for Loop 1; replace by a native TeX diagram in Loop 2.
% Authority: Exp3 local p. 52; combined reader p. 162.
\begin{center}
 \SGAThreeNativeInput[0.96\textwidth]
   {figures/exp3/diagram_039_exp3_p052_remark480_exact_rows.tex}
\end{center}

由蛇引理可得\footnote{编者注: 原版\SGARefLink{sga3:iii:paragraph:4-10}{注 4.10} 中陈述此结果时,
附加了 \(Y'\) 在 \(X'\) 中为局部完全交的假设.
因此, 这一假设也出现在\SGARefLink{sga3:iii:lemma:4-12}{4.12}--\SGARefLink{sga3:iii:corollary:4-14}{4.14} 的陈述中;
事实上它似乎是多余的, 故已从这些陈述中删除.}
\begin{equation}
 \mathscr N_{Y/X}\otimes_{\mathscr O_Y}\mathscr O_{Y'}
 \xrightarrow{\sim}\mathscr N_{Y'/X'}
 \qquad\text{若 \(Y\) 在 \(S\) 上平坦.}
 \tag{4.8.0}\label{sga3:iii:remark:4-8-0:equation:4-8-0}
\end{equation}

\medskip
\noindent\SGAInlineRefTarget{sga3:iii:proposition:4-8}\textbf{命题 4.8.}
\textit{设 \(S,S_0,S_{\mathscr J}\) 及
\(\mathscr I,\mathscr J\) 如\SGARefLink{sga3:iii:proposition:4-5}{4.5} 所述.}\footnote{编者注: 已删除 \(\mathscr I\) 幂零的假设,
因为它似乎是多余的 (参见证明).}
\textit{设 \(X\) 为 \(S\)-概形, \(Y\) 为 \(X\) 的子概形,
\(i:Y\hookrightarrow X\) 为该浸入.}
\begin{enumerate}
\item[\SGAInlineRefTarget{sga3:iii:proposition:4-8:item:i}(i)] \textit{对每个满足
\(f_{\mathscr J}:T_{\mathscr J}\to X_{\mathscr J}\) 经由
\(Y_{\mathscr J}\) 分解的 \(S\)-态射 \(f:T\to X\), 可定义障碍}
  \begin{equation}
   c(X,Y,f)\in
   \operatorname{Hom}_{\mathscr O_{T_0}}
   \left(
    f_0^*\bigl(\mathscr N_{Y/X}
      \otimes_{\mathscr O_Y}\mathscr O_{Y_0}\bigr),
    \mathscr J\mathscr O_T
   \right)
   \tag{*}\label{sga3:iii:proposition:4-8:equation:star}
  \end{equation}
  \textit{它消失当且仅当 \(f\) 存在经由 \(Y\) 的分解.}

\item[\SGAInlineRefTarget{sga3:iii:proposition:4-8:item:ii}(ii)] \textit{设 \(Y'\) 为 \(X\) 的另一个子概形. 假设
\(Y'_{\mathscr J}=Y_{\mathscr J}\), 且 \(Y\) 和 \(Y'\) 在 \(S\) 上平坦.
则存在同构 (参见\SGARefLink{sga3:iii:remark:4-8-0}{4.8.0})}
  \[
   \mathscr J\mathscr O_Y
   \simeq
   \mathscr J\otimes_{\mathscr O_{S_0}}\mathscr O_{Y_0}
   \simeq
   \mathscr J\mathscr O_{Y'},
   \qquad
   \mathscr N_{Y/X}\otimes_{\mathscr O_Y}
     \mathscr O_{Y_{\mathscr J}}
   \xrightarrow{\sim}
   \mathscr N_{Y_{\mathscr J}/X_{\mathscr J}},
  \]
  \textit{从而存在同构}
  \[
   \begin{split}
   u:\ &\operatorname{Hom}_{\mathscr O_{Y_0}}
   \left(
    \mathscr N_{Y_{\mathscr J}/X_{\mathscr J}}
      \otimes_{\mathscr O_{Y_{\mathscr J}}}\mathscr O_{Y_0},
    \mathscr J\mathscr O_Y
   \right)\\
   &\xrightarrow{\sim}
   \operatorname{Hom}_{\mathscr O_{Y_0}}
   \left(
    \mathscr N_{Y/X}\otimes_{\mathscr O_Y}\mathscr O_{Y_0},
    \mathscr J\mathscr O_{Y'}
   \right).
   \end{split}
  \]
  \textit{用 \(i':Y'\hookrightarrow X\) 表示典范浸入,
用 \(d(Y,Y')\) 表示\SGARefLink{sga3:iii:remark:4-5-1}{4.5.1} 中的偏差, 则有}\footnote{编者注: 亦见下文\SGARefLink{sga3:iii:lemma:4-27}{4.27}.}
  \begin{equation}
   c(X,Y,i')=u\bigl(d(Y,Y')\bigr).
   \tag{**}\label{sga3:iii:proposition:4-8:equation:double-star}
  \end{equation}

\item[\SGAInlineRefTarget{sga3:iii:proposition:4-8:item:iii}(iii)] \textit{典范态射}
  \[
   \mathscr N_{Y/X}\xrightarrow{D}i^*(\Omega^1_{X/S})
  \]
  \textit{(参见 SGA 1, II, 公式 4.3)}\footnote{编者注: 亦见 EGA IV\textsubscript{4}, 16.4.21.
回顾, 若 \(U\) 为 \(X\) 的仿射开子集, 使得 \(Y\cap U\) 由
\(A=\mathscr O_X(U)\) 的理想 \(I\) 定义,
\(d\) 表示微分 \(A\to\Gamma(U,\Omega^1_{X/S})\), 且 \(x\in I\),
则 \(D(x+I^2)\) 是 \(\Gamma(U,\Omega^1_{X/S})\otimes_A(A/I)\)
的元素 \(d(x)\otimes1\).}
  \textit{诱导态射}
  \[
   D_0:
   \mathscr N_{Y/X}\otimes_{\mathscr O_Y}\mathscr O_{Y_0}
   \longrightarrow
   \Omega^1_{X_0/S_0}\otimes_{\mathscr O_{X_0}}\mathscr O_{Y_0},
  \]
  \textit{从而对\SGARefLink{sga3:iii:proposition:4-8:item:i}{(i)} 中的每个
\(S\)-态射 \(f:T\to X\), 诱导态射}
  \[
   \begin{split}
   v_{f_0}:\ &
   \operatorname{Hom}_{\mathscr O_{T_0}}
   \left(f_0^*(\Omega^1_{X_0/S_0}),\mathscr J\mathscr O_T\right)
   \\[-0.15em]
   &\longrightarrow
   \operatorname{Hom}_{\mathscr O_{T_0}}
   \left(
    f_0^*\bigl(\mathscr N_{Y/X}
      \otimes_{\mathscr O_Y}\mathscr O_{Y_0}\bigr),
    \mathscr J\mathscr O_T
   \right),\\
   &\hspace{8em}a\longmapsto a\circ f_0^*(D_0),
   \end{split}
  \]
  \textit{其中上面的第一个群是
\(\operatorname{Hom}_{X^+}(T,L_X)\), 参见\SGARefLink{sga3:iii:definition:0-1-5}{0.1.5}.
对 \(a\in\operatorname{Hom}_{X^+}(T,L_X)\), 有}
  \begin{equation}
   c(X,Y,a\cdot f)-c(X,Y,f)=v_{f_0}(a),
   \tag{***}\label{sga3:iii:proposition:4-8:equation:double-starstar}
  \end{equation}
  \textit{其中 \(a\cdot f\) 表示复合态射}
  \[
   T\xrightarrow{a\times f}L_X\times_{X^+}X\longrightarrow X.
  \]
\end{enumerate}

我们将证明命题的 (i), 把断言 (ii) 和 (iii) 留给读者 (不) 验证;
验证时先归结为仿射情形, 再比较显式定义.\footnote{编者注: 下文补出了这些验证.}

现在证明\SGARefLink{sga3:iii:proposition:4-8:item:i}{(i)}.
态射 \(f:T\to X\) 定义环层态射\footnote{编者注: 一方面, 已删除 \(\mathscr I\) 幂零的假设,
即 \(X_0\) 与 \(X\) 有相同底拓扑空间的假设;
另一方面, 扩充了下面的一句话.}
\[
 \phi:\mathscr O_X\longrightarrow f_*(\mathscr O_T).
\]
设 \(U\) 为 \(X\) 的开子概形, 使 \(Y\) 在其中为闭的.
由于 \(T\) 与 \(T_{\mathscr J}\), 相应地 \(Y_{\mathscr J}\) 与 \(Y\),
具有相同的底空间, \(f\) 的底连续映射将 \(T\) 映入 \(U\).
由于 \(U\) 在 \(X\) 中开, 态射 \(\phi\) 诱导环层态射
\[
 \mathscr O_U=\mathscr O_X|_U\longrightarrow f_*(\mathscr O_T),
\]
即 \(f\) 经由 \(U\) 分解.

% Source page: Exp3 local p. 53; combined p. 163;
% original side-pages 134--135.
因此, 可以只考虑 \(Y\) 为闭的情形, 这时它由理想层 \(\mathscr I_Y\) 定义.
态射 \(f\) 经由 \(Y\) 分解当且仅当复合
\[
 \mathscr I_Y\longrightarrow\mathscr O_X
 \longrightarrow f_*(\mathscr O_T)
\]
为零. 由于 \(f_{\mathscr J}\) 经由
\(Y_{\mathscr J}\) 分解, 复合
\[
 \mathscr I_{Y_{\mathscr J}}\longrightarrow
 \mathscr O_{X_{\mathscr J}}\longrightarrow
 f_*(\mathscr O_{T_{\mathscr J}})
\]
为零. 考虑如下交换图, 其中第一行正合:

% Raster stand-in for Loop 1; replace by a native TeX diagram in Loop 2.
% Authority: Exp3 local p. 53; combined reader p. 163.
\begin{center}
 \SGAThreeNativeInput[0.82\textwidth]
   {figures/exp3/diagram_040_exp3_p053_prop48_obstruction_diagram.tex}
\end{center}

从而 \(\phi\) 将 \(\mathscr I_Y\) 映入
\(f_*(\mathscr J\mathscr O_T)\).\footnote{编者注: 下面的论证作了扩充.} 由于 \(\mathscr J^2=0\),
\(\mathscr O_X\)-模 \(f_*(\mathscr J\mathscr O_T)\) 被
\(\mathscr I_Y\) 零化, 其中 \(\mathscr O_X\)-模结构由 \(\phi\) 诱导.
因此 \(\phi\) 诱导 \(\mathscr O_X\)-模态射
\[
 h:i_*(\mathscr N_{Y/X})
   =\mathscr I_Y/\mathscr I_Y^2
 \longrightarrow f_*(\mathscr J\mathscr O_T).
\]

% Source page: Exp3 local p. 54; combined p. 164;
% original side-page 135.
另一方面, 有如下笛卡尔方图:

% Raster stand-in for Loop 1; replace by a native TeX diagram in Loop 2.
% Authority: Exp3 local p. 54; combined reader p. 164.
\begin{center}
 \SGAThreeNativeInput[0.46\textwidth]
   {figures/exp3/diagram_041_exp3_p054_prop48_cartesian_squares.tex}
\end{center}

这里 \(\tau_{T_0},\tau_{X_0},\tau_{Y_0}\) 是由
\(S_0\hookrightarrow S\) 经基变换得到的闭浸入.
由于 \(\mathscr J\mathscr O_T\) 是被 \(\mathscr I\) 零化的拟凝聚
\(\mathscr O_T\)-模, 存在同构
\[
 \mathscr J\mathscr O_T
 \simeq
 (\tau_{T_0})_*\tau_{T_0}^*(\mathscr J\mathscr O_T),
\]
从而
\[
 f_*(\mathscr J\mathscr O_T)
 \simeq
 (\tau_{X_0})_*(f_0)_*\tau_{T_0}^*(\mathscr J\mathscr O_T).
\]
因此, 由伴随性, \(h\) 对应于 \(\mathscr O_{T_0}\)-模态射
\[
 h_0:
 f_0^*\tau_{X_0}^*i_*(\mathscr N_{Y/X})
 \longrightarrow
 \tau_{T_0}^*(\mathscr J\mathscr O_T).
\]
\smallskip
\noindent{\footnotesize\textit{底本记号说明.} Polo--Gille 版的正文以及此处关于 \(h_0\) 的公式,
将笛卡尔方图和前面的基变换公式中标为 \(\tau_{T_0}\) 的闭浸入印作
\(i_{T_0}\). 我们统一使用 \(\tau_{T_0}\).}

现在
\[
 \tau_{X_0}^*i_*(\mathscr N_{Y/X})
 \simeq
 (i_0)_*\tau_{Y_0}^*(\mathscr N_{Y/X})
 =\mathscr N_{Y/X}\otimes_{\mathscr O_Y}\mathscr O_{Y_0}.
\]
回到通常的记号滥用
\(\tau_{T_0}^*(\mathscr J\mathscr O_T)=\mathscr J\mathscr O_T\),
态射 \(h_0\) 可等同于 \(\mathscr O_{T_0}\)-模态射
\[
 h_0:
 f_0^*\bigl(\mathscr N_{Y/X}
   \otimes_{\mathscr O_Y}\mathscr O_{Y_0}\bigr)
 \longrightarrow\mathscr J\mathscr O_T,
\]
这就是所求障碍 \(c(X,Y,f)\). 这就证明了\SGARefLink{sga3:iii:proposition:4-8:item:i}{(i)}.

当 \(f\) 为浸入 \(i':Y'\hookrightarrow X\) 时, 可见
\(c(X,Y,i')\) 来自同态
\[
 \mathscr I_Y\lhook\joinrel\longrightarrow\mathscr O_X
 \longrightarrow\mathscr O_{Y'},
\]
因而由\SGARefLink{sga3:iii:complement:4-4-1}{4.4.1} 和\SGARefLink{sga3:iii:remark:4-5-1}{4.5.1},
它对应于类 \(d(Y,Y')\). 这就证明了 (ii).

\smallskip
\noindent{\footnotesize\textit{底本记号说明.} 在此处, Polo--Gille 版的证明印作
\(i':Y\hookrightarrow X\), 但\SGARefLink{sga3:iii:proposition:4-8:item:ii}{(ii)}
将 \(i'\) 定义为 \(Y'\) 的浸入. 这里使用后者作为定义域.}

下面证明 (iii). 首先,
\(D:\mathscr N_{Y/X}\to i^*(\Omega^1_{X/S})\) 诱导态射
\[
 \begin{split}
 D_0:\quad
 \tau_{Y_0}^*(\mathscr N_{Y/X})
 &\longrightarrow
 \tau_{Y_0}^*i^*(\Omega^1_{X/S})\\
 &=i_0^*\tau_{X_0}^*(\Omega^1_{X/S}).
 \end{split}
\]
由于 \(X_0=X\times_S S_0\), 有
\[
 \tau_{X_0}^*(\Omega^1_{X/S})\simeq\Omega^1_{X_0/S_0}
 \qquad\text{(参见 EGA IV\textsubscript{4}, 16.4.5).}
\]
这就给出了所述态射
\[
 D_0:
 \mathscr N_{Y/X}\otimes_{\mathscr O_Y}\mathscr O_{Y_0}
 \longrightarrow
 \Omega^1_{X_0/S_0}\otimes_{\mathscr O_{X_0}}\mathscr O_{Y_0}.
\]
在下述注之后, 我们将验证等式 \((***)\).

\medskip
\noindent\SGAInlineRefTarget{sga3:iii:remark:4-9}\textbf{注 4.9.}
障碍 \(c(X,Y,f)\) 可在 \(T\) 上局部计算.
设 \(U=\operatorname{Spec}(C)\) 为 \(T\) 的仿射开子集,
位于 \(X\) 的仿射开子集 \(\operatorname{Spec}(A)\) 之上,
而后者又位于 \(S\) 的仿射开子集 \(\operatorname{Spec}(\Lambda)\) 之上.
设 \(J\subset I\subset\Lambda\), 相应地 \(I_Y\subset A\), 为对应于
\(\mathscr J\subset\mathscr I\), 相应地 \(\mathscr I_Y\) 的理想,
令 \(B=A/I_Y\), 并设 \(\phi:A\to C\) 为对应于 \(f:T\to X\) 的
\(\Lambda\)-代数态射. 由于 \(f(T_{\mathscr J})\subset Y_{\mathscr J}\),
有 \(\phi(I_Y)\subset JC\), 因而 \(\phi\) 诱导 \(\Lambda\)-代数态射
\[
 B\longrightarrow C/JC\longrightarrow C_0=C/IC.
\]
于是障碍 \(c=c(X,Y,f)\) 由下述交换图计算:

% Raster stand-in for Loop 1; replace by a native TeX diagram in Loop 2.
% Authority: Exp3 local p. 54; combined reader p. 164.
\begin{center}
 \SGAThreeNativeInput[0.60\textwidth]
   {figures/exp3/diagram_042_exp3_p054_remark49_local_obstruction.tex}
\end{center}

% Source page: Exp3 local p. 55; combined p. 165;
% original side-page 136.
即在 \(U\) 上, 它是下列集合中满足随后条件的唯一元素:
\[
 \operatorname{Hom}_{C_0}
 \left(I_Y/I_Y^2\otimes_B C_0,JC\right)
 =
 \operatorname{Hom}_{B_0}
 \left(I_Y/I_Y^2\otimes_B B_0,JC\right)
\]
采用显然的记号, 该条件为
\[
 c(\overline x\otimes_B1)=\phi(x)
 \qquad\text{对每个 \(x\in I_Y\).}
\]

\footnote{编者注: 以下内容为新增内容.}
现在可以完成\SGARefLink{sga3:iii:proposition:4-8:item:iii}{4.8(iii)} 的证明.
等式 \((***)\) 可在 \(T\) 上局部验证, 故归结为上述仿射情形.
用 \(d_{A/\Lambda}:A\to\Omega^1_{A/\Lambda}\) 表示微分.
在 \(U\) 上, 元素 \(a\) 对应于下列集合的元素 \(a_U\):
\[
 \begin{split}
 &\operatorname{Hom}_{C_0}
 \left(\Omega^1_{A_0/\Lambda_0}\otimes_{A_0}C_0,JC\right)
 \\[-0.15em]
 &\qquad\simeq
 \operatorname{Hom}_{B_0}
 \left(\Omega^1_{A/\Lambda}\otimes_A B_0,JC\right)
 \simeq
 \operatorname{Hom}_{A}(\Omega^1_{A/\Lambda},JC).
 \end{split}
\]
一方面, \(v_{f_0}(a)\) 在 \(U\) 上对应于
\(a_U\circ D_0\), 其中 \(D_0\) 为如下 \(B\)-模态射:\footnote{编者注: 参见编者注 93.}
\[
 \begin{aligned}
 I_Y/I_Y^2&\longrightarrow\Omega^1_{A/\Lambda}\otimes_A B_0,\\
 x+I_Y^2&\longmapsto d_{A/\Lambda}(x)\otimes1.
 \end{aligned}
\]
另一方面 (参见\SGARefLink{sga3:iii:theorem:0-1-8}{0.1.8} 和\SGARefLink{sga3:iii:corollary:0-1-9}{0.1.9} 的证明),
对应于 \(a\cdot f\) 的 \(\Lambda\)-代数态射 \(\phi':A\to C\)
与 \(\phi\) 之差为由 \(a_U\) 关联的 \(\Lambda\)-导子 \(A\to JC\); 即
\[
 \phi'=\phi+a_U\circ d_{A/\Lambda}
 =\phi+a_U\circ(d_{A/\Lambda}\otimes1).
\]
因此, 记 \(c'=c(X,Y,a\cdot f)\), 对每个 \(x\in I_Y\),
用 \(\overline x\) 表示其在 \(I_Y/I_Y^2\) 中的像, 则有
\[
 \begin{split}
 (c'-c)(\overline x\otimes1)
 &=a_U\bigl(d_{A/\Lambda}(x)\otimes1\bigr)\\
 &=(a_U\circ D_0)(\overline x)
 =v_{f_0}(a)(\overline x).
 \end{split}
\]
这就证明了 \(c'-c=v_{f_0}(a)\).

\medskip
\noindent\SGAInlineRefTarget{sga3:iii:paragraph:4-10}\textbf{4.10.}
原版的\SGARefLink{sga3:iii:paragraph:4-10}{注 4.10} 已删除,
因为加入\SGARefLink{sga3:iii:remark:4-8-0}{注 4.8.0} 后, 它已无必要.

% Continue with 4.11 on Exp3 local p. 55 / combined p. 165 /
% original side-page 136.

% Translation coverage: entire assigned source component, 446 source lines.
% Formula audit: 202 expressions; math unchanged except translated text annotations.
% All source reference keys, labels, footnotes and native diagram calls retained.
